\documentclass[10pt,a4paper]{article}
% common preamble for the WSPR Logger papers
\usepackage[T1]{fontenc}
\usepackage[utf8]{inputenc}
\usepackage{lmodern}
\usepackage{microtype}
\usepackage[a4paper,margin=2.4cm,top=2.6cm,bottom=2.6cm]{geometry}
\usepackage{amsmath}
\usepackage{booktabs}
\usepackage{graphicx}
\usepackage{xcolor}
\usepackage{float}
\usepackage{enumitem}
\usepackage{caption}
\usepackage{fancyhdr}
\usepackage{tabularx}
\usepackage{textcomp}
\usepackage[hidelinks]{hyperref}
\definecolor{accent}{HTML}{2B6CB0}
\captionsetup{font=small,labelfont=bf,labelsep=period,justification=raggedright,singlelinecheck=false}
\setlist{itemsep=2pt,topsep=3pt}
\setlength{\parindent}{0pt}
\setlength{\parskip}{5.5pt plus 1pt}
\renewcommand{\arraystretch}{1.12}
\setcounter{tocdepth}{2}
\pagestyle{fancy}
\fancyhf{}
\renewcommand{\headrulewidth}{0.3pt}
\fancyfoot[C]{\small\thepage}
% Result of `tools/backfill.py --compare 2026-06-01`: six live-logged days rebuilt from wspr.live
\newcommand{\cmpSnaps}{995}
\newcommand{\cmpFrom}{1 June 2026}
\newcommand{\cmpTo}{7 June 2026}
\newcommand{\cmpStored}{41.0}
\newcommand{\cmpRebuilt}{41.0}
\newcommand{\cmpRatio}{1.001}
\newcommand{\cmpMore}{3}
\newcommand{\cmpDxStored}{9.8}
\newcommand{\cmpDxRebuilt}{9.8}
\newcommand{\figwide}[2]{\begin{figure}[htbp]\centering\includegraphics[width=\textwidth]{#1}\caption{#2}\end{figure}}
\newcommand{\figw}[3]{\begin{figure}[htbp]\centering\includegraphics[width=#3\textwidth]{#1}\caption{#2}\end{figure}}
\newcommand{\km}{\,km}
\newcommand{\dB}{\,dB}
\newcommand{\MHz}{\,MHz}
\newcommand{\callsign}[1]{\texttt{#1}}
\newcommand{\keyboxed}[1]{\par\medskip\noindent\fcolorbox{accent}{accent!6}{\begin{minipage}{0.95\textwidth}\small\vspace{3pt}#1\vspace{1pt}\end{minipage}}\par\medskip}
\newcommand{\makeheader}[3]{%
  \begin{center}
    {\Large\bfseries #1\par}\vspace{4pt}
    {\large #2\par}\vspace{8pt}
    {\normalsize Otso Laakso, OH2GAX\par}
    {\small #3\par}
  \end{center}\vspace{2pt}\hrule\vspace{8pt}}

\newcommand{\qcFlagged}{3}
\newcommand{\qcRowsAll}{5\,737}
\newcommand{\qcPairs}{5041}
\newcommand{\qcTol}{500}
\newcommand{\qcCall}{G8LZI/2}
\newcommand{\qcGrid}{JO02gl}
\newcommand{\qcDist}{11\,643}
\newcommand{\qcGridDist}{1\,687}
\newcommand{\qcRows}{3\,483}
\newcommand{\swFrom}{2026-05-08}
\newcommand{\swTo}{2026-11-06}
\newcommand{\swDays}{166}
\newcommand{\swSnap}{18\,885}
\newcommand{\nSeasons}{3}
\newcommand{\nSnap}{54\,868}
\newcommand{\nDays}{495}
\newcommand{\dateFrom}{2024-07-04}
\newcommand{\dateTo}{2026-11-06}
\newcommand{\meanRep}{56.3}
\newcommand{\medRep}{53}
\newcommand{\shareDXsix}{16.3}
\newcommand{\daysDXsix}{360}
\newcommand{\pctDaysDXsix}{73}
\newcommand{\daysDXten}{167}
\newcommand{\pctDaysDXten}{34}
\newcommand{\solarZeroSFI}{76}
\newcommand{\nFofC}{2\,973}
\newcommand{\pctFofC}{16}
\newcommand{\nGaps}{543}
\newcommand{\gapHours}{11528.4}
\newcommand{\nStations}{3\,713}
\newcommand{\nRows}{3\,092\,119}
\newcommand{\nSpots}{42\,723}
\newcommand{\nSpotDays}{490}
\newcommand{\homeShare}{86}
\newcommand{\nTxLoc}{102}
\newcommand{\busyFrom}{7}
\newcommand{\hourMinMed}{2}
\newcommand{\hourMaxMed}{81}
\newcommand{\quietMedMax}{4}
\newcommand{\busyMedAvg}{75}
\newcommand{\qLoMid}{46}
\newcommand{\qHiMid}{99}
\newcommand{\busyTo}{20}
\newcommand{\peakThreeHour}{6}
\newcommand{\peakThreePct}{58}
\newcommand{\lowThreeHour}{13}
\newcommand{\lowThreePct}{22}
\newcommand{\dxPeakHour}{5}
\newcommand{\dxPeakPct}{28}
\newcommand{\dxLowHour}{14}
\newcommand{\dxLowPct}{7}
\newcommand{\dxAtF}{28}
\newcommand{\dxThreeAtF}{46}
\newcommand{\dxAtI}{18}
\newcommand{\dxThreeAtI}{45}
\newcommand{\dxAtBC}{16}
\newcommand{\dxThreeAtBC}{25}
\newcommand{\dxAtBJ}{11}
\newcommand{\dxThreeAtBJ}{54}
\newcommand{\nPairs}{5\,038}
\newcommand{\nBeyondBC}{28}
\newcommand{\nBeyondBCOc}{25}
\newcommand{\nBeyondBCOther}{3}
\newcommand{\qcList}{\callsign{G8LZI\slash 2} (grid JO02gl, listed at 11\,643\km, grid 1\,687\km away); \callsign{SM4HVP} (grid JP70od, listed at 17\,038\km, grid 538\km away); \callsign{SV1CDN} (grid KM17tx, listed at 5\,034\km, grid 2\,535\km away)}
\newcommand{\zneurope}{4207}
\newcommand{\zpeurope}{83.5}
\newcommand{\znnamerica}{120}
\newcommand{\zpnamerica}{2.4}
\newcommand{\znafrica}{60}
\newcommand{\zpafrica}{1.2}
\newcommand{\znpacific}{25}
\newcommand{\zppacific}{0.5}
\newcommand{\znsamerica}{7}
\newcommand{\zpsamerica}{0.1}
\newcommand{\znasia}{9}
\newcommand{\zpasia}{0.2}
\newcommand{\znother}{610}
\newcommand{\zpother}{12.1}
\newcommand{\swPct}{75}
\newcommand{\eastPct}{2}
\newcommand{\northPct}{15}
\newcommand{\secSW}{2865}
\newcommand{\secS}{895}
\newcommand{\secW}{256}
\newcommand{\secSE}{168}
\newcommand{\halfStations}{88}
\newcommand{\topStations}{G4ZFQ, ON5KQ, OE9GHV}
\newcommand{\topStationSnaps}{39\,671}
\newcommand{\snrBandA}{-20}
\newcommand{\snrIqrA}{8}
\newcommand{\snrBandB}{-21}
\newcommand{\snrIqrB}{6}
\newcommand{\snrBandC}{-22}
\newcommand{\snrIqrC}{5}
\newcommand{\snrBandD}{-23}
\newcommand{\snrIqrD}{3}
\newcommand{\snrBandE}{-24}
\newcommand{\snrIqrE}{4}
\newcommand{\snrTopCall}{DK6UG}
\newcommand{\snrTopDB}{+24}
\newcommand{\snrTopDist}{1\,592}
\newcommand{\snrTopDate}{2026-09-17}
\newcommand{\snrLowest}{-35}
\newcommand{\corrDistSNR}{-0.19}
\newcommand{\snrMed}{-21}
\newcommand{\naEarlyH}{5}
\newcommand{\naEarlyPct}{31}
\newcommand{\naDipH}{8}
\newcommand{\naDipPct}{15}
\newcommand{\naLatePct}{9}
\newcommand{\paMornH}{7}
\newcommand{\paMornPct}{11}
\newcommand{\paEveH}{17}
\newcommand{\paEvePct}{12}
\newcommand{\daysNA}{290}
\newcommand{\daysAF}{76}
\newcommand{\daysPA}{161}
\newcommand{\daysSA}{12}
\newcommand{\peaknamerica}{04--07}
\newcommand{\peakpctnamerica}{33}
\newcommand{\peakhnamerica}{10}
\newcommand{\peakafrica}{16--19}
\newcommand{\peakpctafrica}{8}
\newcommand{\peakhafrica}{17}
\newcommand{\peakpacific}{16--19}
\newcommand{\peakpctpacific}{12}
\newcommand{\peakhpacific}{17}
\newcommand{\monBestName}{2026-10}
\newcommand{\monBest}{110}
\newcommand{\monWorstName}{2026-07}
\newcommand{\monWorst}{32}
\newcommand{\monCACEAH}{42}
\newcommand{\monuniqCACEAH}{468}
\newcommand{\monmufCACEAH}{--}
\newcommand{\monsfiCACEAH}{--}
\newcommand{\mondxCACEAH}{9}
\newcommand{\monCACEAI}{62}
\newcommand{\monuniqCACEAI}{660}
\newcommand{\monmufCACEAI}{--}
\newcommand{\monsfiCACEAI}{--}
\newcommand{\mondxCACEAI}{21}
\newcommand{\monCACEAJ}{57}
\newcommand{\monuniqCACEAJ}{457}
\newcommand{\monmufCACEAJ}{--}
\newcommand{\monsfiCACEAJ}{--}
\newcommand{\mondxCACEAJ}{16}
\newcommand{\monCACEBA}{61}
\newcommand{\monuniqCACEBA}{445}
\newcommand{\monmufCACEBA}{--}
\newcommand{\monsfiCACEBA}{--}
\newcommand{\mondxCACEBA}{13}
\newcommand{\monCACFAD}{61}
\newcommand{\monuniqCACFAD}{546}
\newcommand{\monmufCACFAD}{--}
\newcommand{\monsfiCACFAD}{--}
\newcommand{\mondxCACFAD}{5}
\newcommand{\monCACFAE}{65}
\newcommand{\monuniqCACFAE}{707}
\newcommand{\monmufCACFAE}{--}
\newcommand{\monsfiCACFAE}{--}
\newcommand{\mondxCACFAE}{13}
\newcommand{\monCACFAF}{73}
\newcommand{\monuniqCACFAF}{704}
\newcommand{\monmufCACFAF}{--}
\newcommand{\monsfiCACFAF}{--}
\newcommand{\mondxCACFAF}{16}
\newcommand{\monCACFAG}{20}
\newcommand{\monuniqCACFAG}{115}
\newcommand{\monmufCACFAG}{--}
\newcommand{\monsfiCACFAG}{--}
\newcommand{\mondxCACFAG}{0}
\newcommand{\monCACFAI}{82}
\newcommand{\monuniqCACFAI}{755}
\newcommand{\monmufCACFAI}{--}
\newcommand{\monsfiCACFAI}{--}
\newcommand{\mondxCACFAI}{22}
\newcommand{\monCACFAJ}{85}
\newcommand{\monuniqCACFAJ}{798}
\newcommand{\monmufCACFAJ}{--}
\newcommand{\monsfiCACFAJ}{--}
\newcommand{\mondxCACFAJ}{17}
\newcommand{\monCACFBA}{103}
\newcommand{\monuniqCACFBA}{773}
\newcommand{\monmufCACFBA}{--}
\newcommand{\monsfiCACFBA}{--}
\newcommand{\mondxCACFBA}{19}
\newcommand{\monCACFBB}{92}
\newcommand{\monuniqCACFBB}{476}
\newcommand{\monmufCACFBB}{--}
\newcommand{\monsfiCACFBB}{--}
\newcommand{\mondxCACFBB}{1}
\newcommand{\monCACGAD}{96}
\newcommand{\monuniqCACGAD}{661}
\newcommand{\monmufCACGAD}{--}
\newcommand{\monsfiCACGAD}{--}
\newcommand{\mondxCACGAD}{31}
\newcommand{\monCACGAE}{90}
\newcommand{\monuniqCACGAE}{723}
\newcommand{\monmufCACGAE}{--}
\newcommand{\monsfiCACGAE}{--}
\newcommand{\mondxCACGAE}{25}
\newcommand{\monCACGAF}{85}
\newcommand{\monuniqCACGAF}{745}
\newcommand{\monmufCACGAF}{19.8}
\newcommand{\monsfiCACGAF}{121}
\newcommand{\mondxCACGAF}{20}
\newcommand{\monCACGAG}{49}
\newcommand{\monuniqCACGAG}{599}
\newcommand{\monmufCACGAG}{18.8}
\newcommand{\monsfiCACGAG}{139}
\newcommand{\mondxCACGAG}{16}
\newcommand{\monCACGAH}{32}
\newcommand{\monuniqCACGAH}{514}
\newcommand{\monmufCACGAH}{18.1}
\newcommand{\monsfiCACGAH}{133}
\newcommand{\mondxCACGAH}{16}
\newcommand{\monCACGAI}{65}
\newcommand{\monuniqCACGAI}{642}
\newcommand{\monmufCACGAI}{19.9}
\newcommand{\monsfiCACGAI}{117}
\newcommand{\mondxCACGAI}{13}
\newcommand{\monCACGAJ}{99}
\newcommand{\monuniqCACGAJ}{707}
\newcommand{\monmufCACGAJ}{22.5}
\newcommand{\monsfiCACGAJ}{106}
\newcommand{\mondxCACGAJ}{14}
\newcommand{\monCACGBA}{110}
\newcommand{\monuniqCACGBA}{765}
\newcommand{\monmufCACGBA}{24.4}
\newcommand{\monsfiCACGBA}{110}
\newcommand{\mondxCACGBA}{14}
\newcommand{\monCACGBB}{103}
\newcommand{\monuniqCACGBB}{335}
\newcommand{\monmufCACGBB}{26.8}
\newcommand{\monsfiCACGBB}{101}
\newcommand{\mondxCACGBB}{5}
\newcommand{\nFullSw}{6}
\newcommand{\corrMonSFI}{-0.88}
\newcommand{\corrMonMUF}{0.92}
\newcommand{\yrCACEDays}{82}
\newcommand{\yrCACEUniq}{1\,098}
\newcommand{\yrCACEMed}{55}
\newcommand{\yrCACEDx}{15.5}
\newcommand{\yrCACEDxDays}{71}
\newcommand{\yrCACEMax}{15\,619}
\newcommand{\yrCACESnap}{8\,730}
\newcommand{\yrCACFDays}{188}
\newcommand{\yrCACFUniq}{2\,045}
\newcommand{\yrCACFMed}{79}
\newcommand{\yrCACFDx}{14.6}
\newcommand{\yrCACFDxDays}{66}
\newcommand{\yrCACFMax}{15\,619}
\newcommand{\yrCACFSnap}{21\,424}
\newcommand{\yrCACGDays}{225}
\newcommand{\yrCACGUniq}{2\,112}
\newcommand{\yrCACGMed}{82}
\newcommand{\yrCACGDx}{18.1}
\newcommand{\yrCACGDxDays}{79}
\newcommand{\yrCACGMax}{15\,619}
\newcommand{\yrCACGSnap}{24\,714}
\newcommand{\nCommon}{3}
\newcommand{\commonNames}{August, September, October}
\newcommand{\coCACEMed}{60}
\newcommand{\coCACEDx}{17.9}
\newcommand{\coCACEDays}{58}
\newcommand{\coCACFMed}{89}
\newcommand{\coCACFDx}{19.3}
\newcommand{\coCACFDays}{88}
\newcommand{\coCACGMed}{96}
\newcommand{\coCACGDx}{13.6}
\newcommand{\coCACGDays}{82}
\newcommand{\nMonCACE}{4}
\newcommand{\nMonCACF}{7}
\newcommand{\nMonCACG}{8}
\newcommand{\kIdxZero}{113}
\newcommand{\kIdxFive}{88}
\newcommand{\kDxZero}{22}
\newcommand{\kDxOne}{19}
\newcommand{\kDxHi}{7}
\newcommand{\aDxLow}{20}
\newcommand{\aDxHigh}{2.7}
\newcommand{\aIdxHigh}{66}
\newcommand{\aIdxLow}{111}
\newcommand{\mufIdxLow}{27}
\newcommand{\mufIdxMid}{144}
\newcommand{\mufDxLow}{4}
\newcommand{\mufDxMid}{19}
\newcommand{\rMufRaw}{0.70}
\newcommand{\rMufRawLo}{0.67}
\newcommand{\rMufRawHi}{0.74}
\newcommand{\rMufRel}{0.43}
\newcommand{\rMufRelLo}{0.37}
\newcommand{\rMufRelHi}{0.48}
\newcommand{\rSfiRel}{0.05}
\newcommand{\rSfiRelLo}{-0.05}
\newcommand{\rSfiRelHi}{0.15}
\newcommand{\rSfiRaw}{-0.13}
\newcommand{\rKRel}{-0.16}
\newcommand{\rKRelLo}{-0.22}
\newcommand{\rKRelHi}{-0.10}
\newcommand{\caseWorstDay}{2026-07-04}
\newcommand{\caseWorstRep}{10}
\newcommand{\caseJulyK}{7}
\newcommand{\caseJulyA}{74}
\newcommand{\caseJulyRep}{10}
\newcommand{\caseTypRep}{80}
\newcommand{\sfiMax}{203}
\newcommand{\worstAllDay}{2025-03-25}
\newcommand{\worstAllRep}{8}
\newcommand{\worstSwDay}{2026-07-04}
\newcommand{\worstSwRep}{10}
\newcommand{\sfiMin}{89}
\newcommand{\covDays}{381}
\newcommand{\covPct}{77}
\newcommand{\covDaysSw}{131}
\newcommand{\covPctSw}{79}
\newcommand{\recBest}{15\,619}
\newcommand{\recBestCall}{VK7JJ}
\newcommand{\recBestDate}{2024-07-10}
\newcommand{\nOceania}{25}
\newcommand{\cmpDays}{164}
\newcommand{\underDays}{68}
\newcommand{\underPct}{41}
\newcommand{\underExDay}{2026-10-07}
\newcommand{\underExStored}{4\,553}
\newcommand{\underExReplay}{15\,614}

\providecommand{\wFixes}{15\,942}
\providecommand{\wFirst}{2026-05-03}
\providecommand{\wLast}{2026-11-05}
\providecommand{\wSpan}{187}
\providecommand{\wDaysWith}{168}
\providecommand{\wLocs}{72}
\providecommand{\wSqNS}{4.6}
\providecommand{\wSqEW}{4.6}
\providecommand{\wRMS}{1.9}
\providecommand{\wHalfDiag}{3.3}
\providecommand{\wSqArea}{21.3}
\providecommand{\wBits}{24.2}
\providecommand{\wNSq}{18\,662\,400}
\providecommand{\wBitsFour}{15.0}
\providecommand{\wSqFourNS}{111}
\providecommand{\wSqFourEW}{110}
\providecommand{\wMinuteOnly}{2}
\providecommand{\wGapBA}{94}
\providecommand{\wGapDA}{98}
\providecommand{\wGapGA}{99}
\providecommand{\wGapMed}{10}
\providecommand{\wGapGh}{75}
\providecommand{\wPerDayMed}{98}
\providecommand{\wPerDayMax}{143}
\providecommand{\wPerDayMean}{95}
\providecommand{\wHeardPct}{68}
\providecommand{\wTxFrac}{18.4}
\providecommand{\wActiveDays}{164}
\providecommand{\wAvSlotAvg}{67}
\providecommand{\wAvDAAvg}{74}
\providecommand{\wAvGAAvg}{77}
\providecommand{\wAvGAMin}{40}
\providecommand{\wAvGAMinH}{2}
\providecommand{\wAvGAMax}{95}
\providecommand{\wAvDAMin}{34}
\providecommand{\wAvDayDA}{90}
\providecommand{\wAvDayGA}{90}
\providecommand{\wAvNightGA}{46}
\providecommand{\wAvSlotDay}{86}
\providecommand{\wAvSlotNight}{31}
\providecommand{\wRepMed}{14}
\providecommand{\wRepMean}{18.5}
\providecommand{\wRepB}{9}
\providecommand{\wRepD}{85}
\providecommand{\wRepBA}{60}
\providecommand{\wRepDMin}{41}
\providecommand{\wRepDMax}{94}
\providecommand{\wMedHMax}{21}
\providecommand{\wHomePct}{84}
\providecommand{\wNear}{11}
\providecommand{\wMid}{2}
\providecommand{\wFar}{3}
\providecommand{\wFarBAA}{3.1}
\providecommand{\wMaxDist}{193}
\providecommand{\wMaxLoc}{KP31kq}
\providecommand{\wMaxDate}{2026-06-19}
\providecommand{\wKPCAkg}{11}
\providecommand{\wTrips}{5}
\providecommand{\wTripDays}{8}
\providecommand{\wMoves}{296}
\providecommand{\wMoveMed}{4.6}
\providecommand{\wMovePJA}{15}
\providecommand{\wMoveMax}{23}
\providecommand{\wSpeedPJA}{87}
\providecommand{\wChangePct}{2}
\providecommand{\wSteps}{14\,985}
\providecommand{\wHomeWk}{61}
\providecommand{\wHomeWe}{88}
\providecommand{\wNAway}{529}
\providecommand{\wPairDays}{4}
\providecommand{\wRelAway}{153}
\providecommand{\wRelAwayLo}{87}
\providecommand{\wRelAwayHi}{338}
\providecommand{\wRelHome}{100}
\providecommand{\wFloor}{-34}
\providecommand{\wSnrMin}{-35}
\providecommand{\wSnrMax}{24}
\providecommand{\wSnrMed}{-19}
\providecommand{\wSnrQB}{-24}
\providecommand{\wSnrQD}{-14}
\providecommand{\wSnrPF}{-28}
\providecommand{\pwTwoH}{100}
\providecommand{\pwDxTwoH}{100}
\providecommand{\pwRepTwoH}{100}
\providecommand{\pwOneH}{100}
\providecommand{\pwDxOneH}{99}
\providecommand{\pwRepOneH}{99}
\providecommand{\pwFifty}{97}
\providecommand{\pwDxFifty}{93}
\providecommand{\pwRepFifty}{94}
\providecommand{\pwTwentyFive}{88}
\providecommand{\pwDxTwentyFive}{74}
\providecommand{\pwRepTwentyFive}{80}
\providecommand{\pwTen}{72}
\providecommand{\pwDxTen}{46}
\providecommand{\pwRepTen}{57}
\providecommand{\pwFive}{58}
\providecommand{\pwDxFive}{25}
\providecommand{\pwRepFive}{40}
\providecommand{\pwOne}{30}
\providecommand{\pwDxOne}{7}
\providecommand{\pwRepOne}{14}
\providecommand{\wKmW}{78\,095}
\providecommand{\wRecord}{15\,619}
\providecommand{\wAbove}{57}
\providecommand{\wStrongMargin}{15}
\providecommand{\eCapWh}{111}
\providecommand{\eAvgmA}{128}
\providecommand{\eAvgW}{0.64}
\providecommand{\eDays}{6.2}
\providecommand{\eUsable}{94}
\providecommand{\eDaily}{15.3}
\providecommand{\eIdleOnlyDays}{7.9}
\providecommand{\eTxOnlyDays}{3.1}
\providecommand{\eObsLo}{4}
\providecommand{\eObsHi}{5}
\providecommand{\eImpliedLo}{65}
\providecommand{\eImpliedHi}{81}

\fancyhead[L]{\small WSPR propagation study --- OH2GAX}
\fancyhead[R]{\small v1.2, 6 November 2026}
\hypersetup{pdftitle={Twenty-metre WSPR propagation observed from a mobile 200 mW beacon},pdfauthor={Otso Laakso, OH2GAX}}
\begin{document}
\makeheader{Twenty-Metre WSPR Propagation Observed from a Mobile 200\,mW Beacon in Southern Finland}{Diurnal, geographic, seasonal and space-weather patterns over three summer seasons, \dateFrom\ to \dateTo}{WSPR Logger project note --- version 1.1 --- 6 November 2026 --- \texttt{github.com/oh2gax/wspr\_logger}}

\begin{abstract}\noindent
A 200\,mW WSPR beacon on 14\,MHz (callsign OH2GAX, Diamond monoband whip on a car roof, home locator KP20jg) has been logged by the WSPR Logger software. This note analyses three summer seasons (2024, 2025 and 2026): \nDays\ calendar days with data between \dateFrom\ and \dateTo, \nSnap\ poll snapshots and \nRows\ individual reception reports from \nStations\ distinct receiving stations. The first two seasons were rebuilt from the public \texttt{wspr.live} database; the solar and ionospheric indices (SFI, K and A index, X-ray class, Juliusruh MUF and foF2) exist only for the \swDays\ days logged live in 2026 (\swFrom\ to \swTo), and the space-weather analysis is confined to them. The band is effectively closed between about 00 and 04~UTC (median \quietMedMax\ stations or fewer) and carries a median of about \busyMedAvg\ stations between 08 and 16~UTC. \zpeurope\,\% of the receiving stations are in Europe and \swPct\,\% lie between south and west-south-west of the transmitter. At least one station farther than 6\,000\km\ was heard on \daysDXsix\ of \nDays\ days (\pctDaysDXsix\,\%), and one farther than 10\,000\km\ on \daysDXten\ days; the record is \recBest\km\ (\recBestCall, \recBestDate), i.e.\ \wKmW\,km per watt. Long-distance windows are concentrated to the early morning and late morning (North America, \peaknamerica\ UTC, with a second maximum near \peakhnamerica~UTC) and late afternoon (Africa and Oceania, \peakafrica\ UTC). Daytime reporting is lowest in midsummer and highest in autumn: in 2026 the median fell from \monCACGAF\ stations in May to \monCACGAH\ in July and rose to \monCACGBA\ in October. In the months observed in all three seasons (\commonNames) the median daytime count was \coCACEMed, \coCACFMed and \coCACGMed\ stations in 2024, 2025 and 2026, and the share of snapshots with a station beyond 6\,000\km\ was \coCACEDx, \coCACFDx\ and \coCACGDx\,\%. In the live-logged 2026 data, after removing the diurnal rhythm, reporting is strongly tied to the Juliusruh MUF (index \mufIdxLow\ for MUF $<$ 12\MHz\ versus \mufIdxMid\ at 20--24\MHz) and weakly and negatively to the geomagnetic K-index (correlation \rKRel); the 10.7\,cm solar flux alone explains nothing ($r=\rSfiRel$). The worst day of the live-logged data (4~July 2026, K\,=\,\caseJulyK, A\,=\,\caseJulyA, \caseJulyRep\ stations against a typical \caseTypRep) followed an M-class flare period with SFI above 195.
\end{abstract}
\vspace{-4pt}
{\small\tableofcontents}


\section{Introduction}
The Weak Signal Propagation Reporter (WSPR) protocol turns every amateur transmitter into a propagation probe: a two-minute, 6\,Hz-wide signal is decoded by volunteer receivers around the world and each decode is stored with the receiver position, signal-to-noise ratio (SNR) and distance~\cite{wsprnet,wsjtx}. A 200\,mW beacon is a modest probe, but because the same transmitter transmits on the same schedule every day, the stream of receptions is a long, homogeneous time series of the 20\,m band as seen from one place.

The WSPR Logger started as a live dashboard for such a beacon: it polls the public WSPR database \texttt{wspr.live}~\cite{wsprlive} every ten minutes, draws the reporters on a map, keeps a log of the farthest reporter per transmission and records the solar conditions at the time. A second database, the \emph{replay database}, stores every poll snapshot in full. The present note is the first systematic analysis of that archive, extended by a rebuild of the two earlier seasons from \texttt{wspr.live}. It asks four questions:
\begin{enumerate}[label=(\alph*)]
\item When is the band open, and when are long-distance paths available?
\item In which directions and at which distances does a mobile 200\,mW beacon in Finland actually get heard?
\item How does reporting change over the season?
\item How much of the variation is explained by the space-weather indices the logger records?
\end{enumerate}
The note is deliberately descriptive. Where a number can be given with a confidence interval it is; where the design of the data collection limits what can be concluded, this is stated (Section~\ref{sec:limits}).

\section{Station, data and definitions}
\subsection{Transmitting station}
The transmitter is a ZachTek WSPR desktop transmitter~\cite{zachtek} with built-in GPS (used for timing and for the position), producing about 200\,mW (23\,dBm) into 50\,$\Omega$. The antenna is a Diamond monoband 14\,MHz whip on a magnetic mount on the roof of a VW Touran. The home locator is KP20jg (60.27\textdegree N, 24.79\textdegree E), but the car moves; \homeShare\,\% of the \nSpots\ logged transmissions carry the home locator and \nTxLoc\ different six-character locators occur in total (the mobile aspect is the subject of a companion paper on position reporting). One WSPR transmission is made every ten minutes.

\subsection{Reception data}\label{sec:data}
Every ten minutes (at :08, :18, \dots) the logger queries \texttt{wspr.live} for the receptions of the last 45 minutes and stores, for the current \emph{snapshot}, every station that has reported the beacon within the preceding 60 minutes: callsign, Maidenhead grid, SNR (dB, 2\,500\,Hz reference bandwidth) and great-circle distance to the transmitter. The live replay archive starts on \swFrom. The two earlier seasons (from 1 July 2024 to 7 May 2026) were \emph{rebuilt}: the script \texttt{tools/backfill.py} downloads every reception of the callsign on 20\,m from \texttt{wspr.live} and reconstructs, for each poll time, the list of stations heard in the preceding 60 minutes. Winter months in which the transmitter was not used simply contain no snapshots. The result is \nSnap\ snapshots on \nDays\ calendar days (\yrCACEDays, \yrCACFDays\ and \yrCACGDays\ days in 2024, 2025 and 2026); within the days with data the snapshots cover about \covPct\,\% of the 24 hours, the rest being times with no report or with the logger offline. The replay database holds \nRows\ rows. A second, older database stores one record per heard transmission (\nSpots\ records on \nSpotDays\ days) with the number of reporters and the largest distance.

The rebuilt snapshots are not guaranteed to equal what the live poll would have stored. The live poll sees only the reports that had reached \texttt{wspr.live} when it ran, and distant stations tend to upload later (Section~\ref{sec:limits}); the rebuild, made afterwards, also sees those late reports, so rebuilt snapshots could in principle list more stations, especially distant ones. This was measured: the live-logged days \cmpFrom\ to \cmpTo\ were rebuilt from \texttt{wspr.live} and compared with the stored snapshots. \cmpSnaps\ snapshots matched (one stored snapshot had no rebuilt counterpart, none the other way); the mean number of stations per snapshot was \cmpStored\ stored against \cmpRebuilt\ rebuilt (ratio \cmpRatio), the rebuild listed more stations in \cmpMore\,\% of the snapshots and never fewer, and a station beyond 6\,000\km\ was present in \cmpDxStored\,\% of the snapshots in both. The effect of late reports is therefore negligible for the statistics below, and differences between seasons are not an artefact of the rebuild (the check covers one week of one season, so it is evidence, not proof, for the others).

\subsection{Solar-terrestrial data}
At every live poll the logger stores the 10.7\,cm solar flux (SFI), the planetary K and A indices, the GOES X-ray class, and the maximum usable frequency (MUF, 3\,000\km\ reference) and critical frequency foF2 of the Juliusruh ionosonde (54.6\textdegree N, 13.4\textdegree E)~\cite{giro,juliusruh}. Indices come from NOAA SWPC-derived feeds~\cite{swpc,hamqsl}. These readings exist for the \swSnap\ snapshots of the \swDays\ live-logged days (\swFrom\ to \swTo); the rebuilt seasons carry none, so everything that involves space weather (Section~\ref{sec:sw} and the case study) is limited to 2026. Readings of zero in SFI or MUF are feed failures (\solarZeroSFI\ snapshots) and are treated as missing. foF2 is available for only \pctFofC\,\% of those snapshots (\nFofC) because the earlier foF2 source was unreliable; it is not used quantitatively here.

\subsection{Definitions}
\begin{description}[leftmargin=1.2em,style=nextline]
\item[Reporters $n_t$] number of distinct stations that reported the beacon in the 60 minutes before snapshot $t$.
\item[Expected reporters $\bar n_h$] mean of $n_t$ over all snapshots in UTC hour $h$ (the typical value for that hour); in the space-weather analysis the mean is taken over the live-logged snapshots only.
\item[Index] $100\sum n_t/\sum\bar n_{h(t)}$ over a set of snapshots. A value of 100 means ``typical for that hour of the day''; this removes the strong day--night rhythm, which would otherwise masquerade as a solar effect.
\item[DX window] snapshot in which at least one reporter was farther than 6\,000\km. DX share is the percentage of snapshots that are DX windows.
\item[Daytime] 08--16~UTC (10--18 or 11--19 local time, depending on daylight saving).
\end{description}

\subsection{Quality control}\label{sec:qc}
Each (callsign, grid) pair was compared with the distance the logger stored: a pair is rejected if the distance computed from the receiver's grid differs from the reported distance by more than \qcTol\km. Of \qcPairs\ pairs, \qcFlagged\ were rejected: \qcList. Their \qcRowsAll\ rows were dropped from every statistic below. Without this step the farthest ``reporter'' of the whole record would have been spurious. All station and distance results use the cleaned data.

\section{Methods}
\textbf{Confidence intervals.} Proportions of snapshots (DX shares) carry 95\,\% Wilson intervals. Indices and correlations carry 95\,\% day-block bootstrap intervals (2\,000 resamples of whole days with replacement), because consecutive snapshots share 50 of their 60 minutes and are strongly autocorrelated; treating them as independent would make every interval far too narrow.

\textbf{Regions.} Receiving stations were assigned to a region from the centre of their four-character grid. \textbf{Sectors.} Bearing from the transmitter home locator is divided into eight 45\textdegree\ sectors. \textbf{Distance bands.} 0--1\,500, 1\,500--3\,000, 3\,000--6\,000, 6\,000--10\,000 and $>$10\,000\km. A \emph{station} statistic uses the median over the snapshots in which that station appears.

\textbf{Space-weather bins.} MUF, K and A are binned and, within each bin, the index and DX share are computed over all snapshots in the bin. Bins with fewer than 30 snapshots are omitted. Correlations use the per-snapshot index $n_t/\bar n_{h(t)}$ (relative reporting) and the raw $n_t$ for comparison.

\textbf{Reproducibility.} All figures and numbers in this note are generated by \texttt{analyze.py} from the two databases; the \LaTeX\ macros in the text are written by the script, so no value is typed by hand (Appendix~\ref{app:repro}).

\section{Results}
\subsection{Overview}
The record contains \nSnap\ snapshots from \dateFrom\ to \dateTo, with a mean of \meanRep\ and a median of \medRep\ reporters per snapshot; \nStations\ distinct callsigns appear (\nPairs\ callsign--grid pairs). \pctDaysDXsix\,\% of days (\daysDXsix\ of \nDays) had at least one reporter beyond 6\,000\km, \daysDXten\ days ( \pctDaysDXten\,\%) beyond 10\,000\km. \shareDXsix\,\% of all snapshots are DX windows. Figure~\ref{fig:heat} shows the structure at a glance: three seasons separated by the winter break, a regular daytime band of strong reporting, a dead band in the small hours, occasional dropouts of the logger, a weaker midsummer, and a brightening in September--October. The seasons are compared in Section~\ref{sec:seasons}.

\figwide{out/fig_heatmap.pdf}{Band-open heatmap for 20\,m: mean number of stations heard per UTC hour for each calendar day (colour clipped at the 97th percentile). White columns are days without data: the winter break, days on which the transmitter was off, and periods when the logger was not running.\label{fig:heat}}

\subsection{Diurnal pattern}
Figure~\ref{fig:diurnal}a shows the typical day. Between 00 and 04~UTC the median is $\le$\quietMedMax\ stations; reporting rises steeply after 05 and the median is at least 50 stations from \busyFrom\ to \busyTo~UTC; from 08 to 16 UTC it is about \busyMedAvg\ (maximum hourly median \hourMaxMed). The interquartile range is wide, on average \qLoMid\ to \qHiMid\ stations between 08 and 16~UTC, so day-to-day variability is large compared with the typical value.

The probability of a long-distance window is \emph{not} highest when the band is busiest (Figure~\ref{fig:diurnal}b). The share of snapshots with a reporter beyond 6\,000\km\ peaks at \dxPeakPct\,\% around \dxPeakHour~UTC but is only \dxLowPct\,\% at \dxLowHour~UTC; \dxAtF\,\% at 05~UTC and \dxAtBJ\,\% at 19~UTC. A reporter beyond 3\,000\km\ is available in up to \peakThreePct\,\% of snapshots at \peakThreeHour~UTC and in only \lowThreePct\,\% at \lowThreeHour~UTC. A plausible reading is that around noon the F2 layer carries the short and medium paths into Europe while the long paths, which need the layer to be illuminated at both ends or along a grey line, are least likely; the morning and evening maxima coincide with the periods when the path to North America (morning) and to Africa and Oceania (afternoon and evening) passes through twilight. The logger data cannot test this mechanism; they only document the timing.

\figwide{out/fig_diurnal.pdf}{(a) Typical day: median and interquartile range of the stations heard in 60 minutes by UTC hour. (b) Share of snapshots with at least one station beyond 6\,000\km\ (bars, 95\,\% Wilson intervals) and beyond 3\,000\km\ (line).\label{fig:diurnal}}

\subsection{Geography: where the beacon is heard}
Table~\ref{tab:zones} and Figure~\ref{fig:geo} show that the footprint is overwhelmingly European: \zneurope\ of \nPairs\ stations (\zpeurope\,\%). North America contributes \znnamerica\ stations (\zpnamerica\,\%), Africa \znafrica, Oceania \znpacific\ and South America \znsamerica; Asia only \znasia. Measured by bearing, \secSW\ stations lie in the south-west sector (225--270\textdegree), \secS\ in the south sector (180--225\textdegree), \secW\ to the west and \secSE\ to the south-east; the east and north together hold only about \eastPct\,\% and \northPct\,\% of the stations. The receiver network is densest in Central and Western Europe, which are also the directions of the best 20\,m paths from Finland, so this is a joint property of propagation and receiver density, not of propagation alone.

The reporting is also concentrated on a few stations: \halfStations\ stations account for half of all reports (the three most frequent are \topStations; the top station appears in \topStationSnaps\ snapshots). These are mostly continuously running receivers, so the network response to a change in propagation is dominated by a small set of fixed receivers.

\begin{table}[htbp]\centering\small
\caption{Receiving stations by region (unique callsign--grid pairs, quality-controlled).}\label{tab:zones}
\begin{tabular}{lrrrrr}
\toprule
Region of station & Stations & Share (\%) & Median dist. (km) & Max. dist. (km) & Best SNR (dB) \\
\midrule
Europe & 4\,207 & 83.5 & 1\,595 & 4\,682 & +24 \\
North America & 120 & 2.4 & 6\,868 & 9\,102 & +1 \\
Africa & 60 & 1.2 & 3\,160 & 9\,877 & +6 \\
Oceania & 25 & 0.5 & 15\,070 & 15\,619 & +16 \\
South America & 7 & 0.1 & 11\,351 & 13\,421 & -5 \\
Asia & 9 & 0.2 & 2\,470 & 4\,466 & -1 \\
Other / unknown & 610 & 12.1 & 2\,305 & 14\,785 & -5 \\
\bottomrule
\end{tabular}
\end{table}

\figwide{out/fig_geography.pdf}{(a) Position of every receiving station by bearing from the transmitter and distance (dot size grows with the number of snapshots in which the station was heard; colour = region). (b) Median SNR per station by distance band (boxes: quartiles; whiskers: 1.5 interquartile ranges).\label{fig:geo}}

\subsection{Distance and signal strength}
Over all reports the median SNR is \wSnrMed\dB\ with quartiles at \wSnrQB\ and \wSnrQD\dB\ (5th percentile \wSnrPF\dB); the lowest value in the archive is \snrLowest\dB\ and the highest \snrTopDB\dB, measured at \callsign{\snrTopCall} at \snrTopDist\km\ on \snrTopDate\ (Table~\ref{tab:snr}). SNR falls with the logarithm of distance, but weakly (correlation \corrDistSNR\ between $\log_{10}d$ and the station median SNR): the median is \snrBandA\dB\ below 1\,500\km, \snrBandB\ to \snrBandC\dB\ between 1\,500 and 6\,000\km\ and only \snrBandD\ to \snrBandE\dB\ beyond 6\,000\km. The spread narrows with distance (interquartile range \snrIqrA\dB\ below 1\,500\km, \snrIqrD\ to \snrIqrE\dB\ beyond 6\,000\km), and the long-distance stations lie close to the lower decoding limit of the archive, about $-$33\dB, so the long paths operate with little margin and a few dB of extra attenuation removes them. The weak distance dependence is what one expects for a single F2 hop whose loss is dominated by absorption and focusing rather than by free-space spreading; the logger data do not test this explanation.

\begin{table}[htbp]\centering\small
\caption{Station SNR by distance band (median over the snapshots in which a station was heard).}\label{tab:snr}
\begin{tabular}{lrrrrr}
\toprule
Distance (km) & Stations & Median SNR (dB) & Q1 (dB) & Q3 (dB) & Best SNR (dB) \\
\midrule
< 1\,500 & 603 & -19 & -23 & -14 & +15 \\
1\,500--3\,000 & 1\,277 & -22 & -24 & -18 & +24 \\
3\,000--6\,000 & 51 & -22 & -25 & -20 & -5 \\
6\,000--10\,000 & 56 & -24 & -26 & -23 & -7 \\
$>$ 10\,000 & 16 & -25 & -26 & -23 & -5 \\
\bottomrule
\end{tabular}
\end{table}


\subsection{Long-distance windows by region}
Long-distance reports ($>$6\,000\km) come from three regions. North America reported on \daysNA\ days and shows two maxima (Figure~\ref{fig:zones}): one in the early morning, up to \naEarlyPct\,\% of days at \naEarlyH~UTC, and a second in the late morning, \peakpctnamerica\,\% of days at \peakhnamerica~UTC, separated by a dip to \naDipPct\,\% at \naDipH~UTC; from 14~UTC onwards it never exceeds \naLatePct\,\%. Africa reported on \daysAF\ days, concentrated in the afternoon and evening (\peakafrica~UTC; maximum \peakpctafrica\,\% of days at \peakhafrica~UTC). Oceania reported on \daysPA\ days, with a late-afternoon maximum (\paEvePct\,\% at \paEveH~UTC) and a smaller early-morning one (\paMornPct\,\% at \paMornH~UTC). South America contributed on \daysSA\ days only. The farthest of each region is in Table~\ref{tab:dxzones}.

\figwide{out/fig_zonehours.pdf}{Percentage of days with at least one station of the region beyond 6\,000\km\ in each UTC hour, for the three regions that produced long-distance reports.\label{fig:zones}}
\begin{table}[htbp]\centering\small
\caption{Days with at least one station beyond 6\,000 km, by region.}\label{tab:dxzones}
\begin{tabular}{lrrrrl}
\toprule
Region & Days & Best 3-h window (UTC) & Peak share (\%) & Farthest (km) & Station, date \\
\midrule
North America & 96 & 10--13 & 31 & 9\,102 & YV5RAB, 2026-09-24 \\
Africa & 57 & 16--19 & 21 & 9\,877 & ZS3Y, 2026-06-25 \\
Oceania & 53 & 16--19 & 14 & 15\,619 & VK7JJ, 2026-06-23 \\
South America & 1 & -- & -- & 13\,181 & CE3AM, 2026-06-28 \\
\bottomrule
\end{tabular}
\\[2pt]{\footnotesize Share = percentage of days with data in that UTC hour on which at least one station of the region beyond 6\,000 km was heard.}
\end{table}


\subsection{Seasonal variation and the three seasons}
\label{sec:seasons}
The archive covers three seasons of unequal length (Table~\ref{tab:seasons}): July--October 2024 (\yrCACEDays\ days with data), March--November 2025 (\yrCACFDays\ days) and March--November 2026 so far (\yrCACGDays\ days). The transmitter is not used in winter, and in 2025 almost nothing was logged between early June and the end of July. The seasons are therefore compared month by month (Table~\ref{tab:yoy}, Figure~\ref{fig:seasonsfig}).

\textbf{Within a season.} Monthly daytime reporting (Table~\ref{tab:months}, Figure~\ref{fig:months}) shows a pronounced pattern that is \emph{not} the one the solar flux alone would predict. In 2026, the only season with a spring-to-autumn record, the median daytime count was \monCACGAD\ in March, \monCACGAF\ in May, fell to \monCACGAH\ in July and rose to \monCACGBA\ in October. The midsummer trough is also seen in 2024 (\monCACEAH\ in July against \monCACEAI\ in August); in 2025 the midsummer months were not logged, but reporting rose from \monCACFAD\ in March to \monCACFBA\ in October. The mean SFI of 2026 went the other way: \monsfiCACGAF\ in May, \monsfiCACGAG\ in June and \monsfiCACGAH\ in July, down to \monsfiCACGBA\ in October. The correlation of the monthly median reporting with the monthly mean SFI is \corrMonSFI\ and with the monthly median Juliusruh MUF \corrMonMUF\ (\nFullSw\ complete months with solar data, so these are descriptive, not statistical, values).

Two effects probably contribute and cannot be separated in this record. First, the mid-latitude F2 layer has a well-known seasonal anomaly: for the same solar activity the daytime critical frequency is higher in the winter half-year than in summer, which raises the MUF in autumn (here from a daytime median of \monmufCACGAF\MHz\ in May and \monmufCACGAH\MHz\ in July to \monmufCACGBA\MHz\ in October and \monmufCACGBB\MHz\ in November 2026). Second, the receiving network itself changes: the number of distinct stations that heard the beacon fell from \monuniqCACGAF\ in May 2026 to \monuniqCACGAH\ in July and rose to \monuniqCACGBA\ in October, which is consistent with fewer attended and automatic receivers during the northern-hemisphere summer. Without an independent reference beacon, neither explanation can be ruled out.

\textbf{Between seasons.} In the months observed in all three seasons (\commonNames) the median daytime count was \coCACEMed\ stations in 2024, \coCACFMed\ in 2025 and \coCACGMed\ in 2026. The step from 2024 to 2025 coincides with a rise in the number of distinct stations heard in the same months (\monuniqCACEBA\ in October 2024 against \monuniqCACFBA\ in October 2025, \monuniqCACEAI\ against \monuniqCACFAI\ in August); it is probably a mixture of a growing receiver network, a different state of the ionosphere (solar cycle~25 passed its maximum within the period, but the flux of 2024 and 2025 is not archived), the bias of the rebuilt snapshots (Section~\ref{sec:limits}) and possible changes of the station that the data cannot show. August 2026 is the exception to the upward trend (\monCACGAI\ against \monCACFAI\ in 2025). The long-distance share did not follow the daytime count: \coCACEDx\,\% of the snapshots of these months had a station beyond 6\,000\km\ in 2024, \coCACFDx\,\% in 2025 and \coCACGDx\,\% in 2026. Every season produced reports beyond 15\,000\km\ (Table~\ref{tab:seasons}).

\begin{table}[htbp]\centering\small
\caption{The three seasons. Daytime median = median number of stations heard between 08 and 16~UTC; stations = distinct callsigns heard in the season; DX share = snapshots with a station beyond 6\,000\km; days with DX = days with at least one such station. The seasons cover different months (Table~\ref{tab:yoy}).}\label{tab:seasons}
\begin{tabular}{lrrrrrrr}
\toprule
Season & Days & Snapshots & Stations & Daytime median & DX share (\%) & Days with DX (\%) & Farthest (km) \\
\midrule
2024 & 82 & 8\,730 & 1\,098 & 55 & 15.5 & 71 & 15\,619 \\
2025 & 188 & 21\,424 & 2\,045 & 79 & 14.6 & 66 & 15\,619 \\
2026 & 225 & 24\,714 & 2\,112 & 82 & 18.1 & 79 & 15\,619 \\
\bottomrule
\end{tabular}
\end{table}

\begin{table}[htbp]\centering\small
\caption{Month-by-month comparison of the seasons: median number of stations heard between 08 and 16~UTC, with the DX share in per cent in parentheses. Months with fewer than ten days of data are left out (--).}\label{tab:yoy}
\begin{tabular}{lrrr}
\toprule
Month & 2024 & 2025 & 2026 \\
\midrule
Mar & -- & 61 (5) & 96 (31) \\
Apr & -- & 65 (13) & 90 (25) \\
May & -- & 73 (16) & 85 (20) \\
Jun & -- & -- & 49 (16) \\
Jul & 42 (9) & -- & 32 (16) \\
Aug & 62 (21) & 82 (22) & 65 (13) \\
Sep & 57 (16) & 85 (17) & 99 (14) \\
Oct & 61 (13) & 103 (19) & 110 (14) \\
Nov & -- & 92 (1) & -- \\
\bottomrule
\end{tabular}
\end{table}

\figwide{out/fig_seasons.pdf}{The seasons month by month: (a) median number of stations heard between 08 and 16~UTC; (b) share of snapshots with a station beyond 6\,000\km. Only months with at least ten days of data are drawn.\label{fig:seasonsfig}}
\begin{table}[htbp]\centering\small
\caption{Monthly summary. Stations = distinct callsigns heard in the month.}\label{tab:months}
\begin{tabular}{lrrrrrr}
\toprule
Month & Days & Snapshots & Stations & Median stations 08--16 UTC & Median MUF (MHz) & Mean SFI \\
\midrule
2026-05 & 24 & 3\,556 & 665 & 87 & 19.8 & 121 \\
2026-06 & 28 & 3\,762 & 600 & 49 & 18.8 & 139 \\
2026-07 & 26 & 3\,370 & 515 & 32 & 18.1 & 133 \\
2026-08 & 26 & 2\,737 & 642 & 65 & 19.9 & 117 \\
2026-09 & 28 & 2\,597 & 707 & 99 & 22.5 & 106 \\
2026-10 & 28 & 2\,532 & 766 & 110 & 24.4 & 110 \\
2026-11 & 6 & 313 & 331 & 102 & 25.9 & 99 \\
\bottomrule
\end{tabular}
\end{table}

\figwide{out/fig_monthly.pdf}{Monthly median number of stations heard between 08 and 16~UTC (bars, coloured by season) and, for the live-logged months of 2026 only, the median Juliusruh MUF over the same hours (circles) and the mean SFI (squares). November 2026 comprises six days.\label{fig:months}}

\subsection{Space weather}
\label{sec:sw}
Figure~\ref{fig:sw} and Tables~\ref{tab:muf}--\ref{tab:a} relate the hour-adjusted reporting index and the DX share to the three indices recorded at each poll. Only the \swDays\ live-logged days of 2026 (\swFrom\ to \swTo, \swSnap\ snapshots) carry these readings, so this section describes one season.

\textbf{MUF.} The Juliusruh MUF is the strongest single predictor. At MUF below 12\MHz\ the index is only \mufIdxLow\ and the DX share \mufDxLow\,\%; at 20--24\MHz\ the index is \mufIdxMid\ and the DX share \mufDxMid\,\%. The raw correlation between reporters and MUF is \rMufRaw\ (95\,\% CI \rMufRawLo--\rMufRawHi), but much of that is the common diurnal rhythm; after removing it the correlation with relative reporting is \rMufRel\ (\rMufRelLo--\rMufRelHi). The relationship is therefore real but moderate: a high MUF is necessary rather than sufficient.

\textbf{Geomagnetic activity.} Reporting falls with the K-index, but gently: the index is \kIdxZero\ at K\,=\,0 and \kIdxFive\ at K\,$\ge$\,5; the correlation is \rKRel\ (\rKRelLo\ to \rKRelHi). The DX share is far more sensitive: \kDxZero\,\% at K\,=\,0, \kDxOne\,\% at K\,=\,1 and only \kDxHi\,\% on average for K\,$\ge$\,3. The A-index gives the same picture over the day: for A\,$<$\,8 the DX share is \aDxLow\,\% (index \aIdxLow), for A\,$\ge$\,30 it is \aDxHigh\,\% and the index \aIdxHigh. Disturbed conditions thus mostly remove the weakest, longest paths first, which is the expected behaviour of an auroral-zone path at 60\textdegree N.

\textbf{Solar flux.} The solar flux has no detectable effect on hourly reporting on its own ($r=\rSfiRel$, CI \rSfiRelLo\ to \rSfiRelHi; the raw value is \rSfiRaw). This is not a statement that the Sun is irrelevant: SFI changed only between about \sfiMin\ and \sfiMax\ during these days, strong flux comes with flares and storms (Section~\ref{sec:case}), and SFI is a daily number while reporting responds to the hour-to-hour ionosphere.

\begin{table}[htbp]\centering\small
\caption{Reporting versus Juliusruh MUF. Index: stations heard relative to the typical value for the same UTC hour (100 = typical); DX share: snapshots with a station beyond 6\,000 km. Confidence intervals by day-block bootstrap.}\label{tab:muf}
\begin{tabular}{lrrrrrr}
\toprule
MUF (MHz) & Snapshots & Mean stations & Index & 95 \% CI & DX share (\%) & 95 \% CI \\
\midrule
$<$ 12 & 1\,549 & 6 & 27 & 21--33 & 4 & 2--7 \\
12--16 & 3\,494 & 13 & 38 & 34--42 & 10 & 7--13 \\
16--20 & 7\,534 & 52 & 81 & 75--88 & 15 & 13--18 \\
20--24 & 4\,710 & 103 & 144 & 135--153 & 19 & 16--23 \\
24--28 & 1\,160 & 110 & 142 & 130--155 & 20 & 14--27 \\
$\geq$ 28 & 335 & 100 & 126 & 113--138 & 17 & 8--29 \\
\bottomrule
\end{tabular}
\end{table}

\figwide{out/fig_spaceweather.pdf}{Hour-adjusted reporting index (bars, left axis, 100 = typical for the UTC hour; 95\,\% day-block bootstrap intervals) and DX share (green line, right axis) by (a) Juliusruh MUF, (b) K-index and (c) A-index. Blue bars are at or above the typical level, orange below.\label{fig:sw}}
\begin{table}[htbp]\centering\small
\caption{Reporting versus the K-index (day-block bootstrap intervals).}\label{tab:k}
\begin{tabular}{lrrrrrr}
\toprule
K-index & Snapshots & Mean stations & Index & 95 \% CI & DX share (\%) & 95 \% CI \\
\midrule
0 & 2\,325 & 64 & 113 & 102--125 & 22 & 16--28 \\
1 & 6\,238 & 65 & 108 & 98--120 & 19 & 16--22 \\
2 & 5\,850 & 55 & 95 & 87--105 & 12 & 10--14 \\
3 & 3\,174 & 50 & 88 & 75--102 & 9 & 6--13 \\
4 & 880 & 45 & 82 & 59--104 & 6 & 0--12 \\
5+ & 398 & 51 & 88 & 43--143 & 7 & 1--15 \\
\bottomrule
\end{tabular}
\end{table}

\begin{table}[htbp]\centering\small
\caption{Reporting versus the A-index (day-block bootstrap intervals).}\label{tab:a}
\begin{tabular}{lrrrrrr}
\toprule
A-index & Snapshots & Mean stations & Index & 95 \% CI & DX share (\%) & 95 \% CI \\
\midrule
0--7 & 9\,178 & 65 & 111 & 103--122 & 20 & 17--22 \\
8--15 & 6\,200 & 59 & 101 & 88--116 & 12 & 9--15 \\
16--29 & 2\,757 & 39 & 69 & 57--81 & 8 & 5--10 \\
$\geq$ 30 & 732 & 41 & 66 & 34--96 & 3 & 0--7 \\
\bottomrule
\end{tabular}
\end{table}


\subsection{Case study: 24 June -- 9 July}
\label{sec:case}
The most disturbed fortnight of the live-logged data is instructive (Table~\ref{tab:case}, Figure~\ref{fig:case}). Solar flux rose from 131 on 24~June to above 195 on 30~June--3~July with M-class flares on several days. Reporting followed no simple pattern: 2 and 3~July had the best days of the fortnight (about 93 and 70 stations, with K\,=\,2--3) while 25~June (K\,=\,5, A\,=\,26) and 1~July (K\,=\,4, A\,=\,18) were weak. The worst day of the live-logged data was 4~July: K\,=\,\caseJulyK, A\,=\,\caseJulyA, and only \caseJulyRep\ stations in the daytime window against the typical \caseTypRep\ (the lowest daytime mean of the whole archive, \worstAllRep\ stations, is from \worstAllDay, a day without solar data). The disturbance persisted: 5 July (K\,=\,5, A\,=\,74) had 13 stations, and 6~July, with K down to 3 and A\,=\,18, still only 10; reporting recovered to 28 on 7~July and to 67 on 9~July. The pattern, a strong single-day collapse followed by a recovery lag of two to three days, is the classic ionospheric storm signature at high geomagnetic latitudes: it is the K-index and the A-index, not the X-ray class or the flux, that mark the days that failed.

\begin{table}[htbp]\centering\small
\caption{The period 24 June -- 9 July 2026: space weather and reporting (daytime mean).}\label{tab:case}
\begin{tabular}{lrrrrl}
\toprule
Date & SFI & Max K & Max A & Max X-ray & Stations 08--16 UTC \\
\midrule
2026-06-24 & 131 & 3 & 15 & C & 37 \\
2026-06-25 & 139 & 5 & 26 & C & 16 \\
2026-06-26 & 157 & 4 & 26 & C & 24 \\
2026-06-27 & 166 & 3 & 12 & C & 15 \\
2026-06-28 & 188 & 2 & 8 & C & 39 \\
2026-06-29 & 187 & 1 & 8 & C & 27 \\
2026-06-30 & 196 & 4 & 18 & M & 40 \\
2026-07-01 & 203 & 4 & 18 & M & 16 \\
2026-07-02 & 201 & 2 & 17 & M & 93 \\
2026-07-03 & 201 & 3 & 12 & M & 70 \\
2026-07-04 & 184 & 7 & 74 & M & 10 \\
2026-07-05 & 159 & 5 & 74 & M & 13 \\
2026-07-06 & 138 & 3 & 18 & C & 10 \\
2026-07-07 & 123 & 2 & 7 & C & 28 \\
2026-07-09 & 115 & 4 & 16 & C & 67 \\
\bottomrule
\end{tabular}
\end{table}

\figwide{out/fig_case.pdf}{Daytime reporting (bars) with maximum K-index (red) and mean SFI (green) for 24~June -- 9~July 2026 (8~July had no complete daytime data).\label{fig:case}}

\subsection{Records}
All eight farthest stations (Table~\ref{tab:records}) are in Australia; the record is \recBestCall\ at \recBest\km\ on \recBestDate, which corresponds to \wKmW\,km per watt of transmitter output. Three of the eight rows (\callsign{VK7JJ}, \callsign{VK7FLYN}, \callsign{VK7FLYN/2}) share one grid square and are probably one receiving site using several callsigns. Of the \nBeyondBC\ stations that were heard beyond 12\,000\km, \nBeyondBCOc\ are in Oceania, two in South America (Argentina and Chile) and one in Antarctica (\callsign{DP0GVN/1}).

\begin{table}[htbp]\centering\small
\caption{The eight farthest stations (one row per station, quality-controlled).}\label{tab:records}
\begin{tabular}{rlllrl}
\toprule
\# & Station & Region & Grid & Distance (km) & When (UTC) \\
\midrule
1 & VK7JJ & Australia & QE38lr & 15\,619 & 2024-07-10 17:38 \\
2 & VK7FLYN & Australia & QE38lr & 15\,619 & 2026-09-24 06:08 \\
3 & VK7FLYN\slash 2 & Australia & QE38lr & 15\,619 & 2026-09-30 05:48 \\
4 & VK7EKA & Australia & QE38kv & 15\,596 & 2024-10-22 06:28 \\
5 & VK3AG & Australia & QF31 & 15\,385 & 2026-03-08 13:48 \\
6 & SWLVK3 & Australia & QF22xv & 15\,208 & 2025-10-10 06:08 \\
7 & SWL\slash VK3 & Australia & QF12wb & 15\,165 & 2026-06-14 21:38 \\
8 & VK2EFM & Australia & QF56oq & 15\,156 & 2024-08-16 19:48 \\
\bottomrule
\end{tabular}
\end{table}


\section{Discussion}
\label{sec:limits}
\textbf{What the data support.} The diurnal structure, the European concentration and the existence and timing of the North American, African and Oceanian windows are all robust: they rest on hundreds of days or thousands of snapshots and are not sensitive to the quality-control step. The dependence of reporting on MUF and on K/A is also robust in sign and, with day-block intervals, in size, but it rests on one season (2026). The seasonal pattern within a year is real in the data but cannot be attributed, and the differences between the seasons are mixed with changes in the receiver network and with the rebuild of the older snapshots.

\textbf{Limits.}
\begin{itemize}
\item \emph{One path, one band, three summer seasons.} The archive covers July--October 2024, March--November 2025 and March--November 2026 on 14\,MHz only. The car-mounted antenna is not used in winter (snow, low temperature), so no winter data exist and seasonal statements concern March--November. Whether the equipment, the antenna position or the operating habits were identical in the three seasons is not recorded in the data, which limits the comparison between seasons.
\item \emph{Rebuilt snapshots and missing space weather.} The snapshots of 2024, 2025 and early 2026 were reconstructed from \texttt{wspr.live} after the fact. The live poll could not see reports that arrived later, so the rebuilt snapshots could list more stations; a comparison of \cmpSnaps\ live snapshots with their rebuild (Section~\ref{sec:data}) found a ratio of \cmpRatio, i.e.\ no measurable difference, but only for one week in June 2026. No solar or MUF values were rebuilt, so the space-weather results are limited to \swDays\ days.
\item \emph{Receiver network.} Reporting counts depend on how many receivers are listening. The \halfStations\ most frequent stations account for half of all reports, so a station going off-air shifts the counts without any change in propagation. The sparsely covered directions (north, east) say nothing about propagation there.
\item \emph{Mobile transmitter.} The car moves; \homeShare\,\% of the logged transmissions come from the home square. Position-dependent effects (antenna surroundings, local noise at the receiving end of a path) are not separated. The companion paper on position reporting examines the movement itself.
\item \emph{Overlapping windows.} Consecutive snapshots overlap by 50 minutes; all intervals use day-block resampling for this reason, but the effective sample size is closer to the number of days than to the number of snapshots.
\item \emph{Indices versus path.} Juliusruh is about 900\km\ south-west of the transmitter and not on the path to North America; K and A are planetary values, and SFI is daily. The indices are proxies.
\item \emph{SNR.} Reported SNR is estimated by the receiving station and depends on its noise floor and software; only medians over many stations are interpreted.
\end{itemize}

\section{Data quality and software notes}
Two issues were found while preparing the study and are relevant to anyone reading the archive.

\textbf{Under-reported Best DX.} The logger's per-transmission record (used for the ``Best DX'' figure on the live page) was written from a single poll six minutes after the transmission. Distant stations upload to WSPR~Net later, so the farthest reporter of a transmission was often missing; comparing the two databases day by day over the live-logged days showed that on \underDays\ of \cmpDays\ days (\underPct\,\%) the farthest reporter in the replay data was farther than the stored value (largest difference on \underExDay: \underExStored\km\ stored against \underExReplay\km). The poll now covers the last 45 minutes and raises the stored value as later reports arrive; the backfill script raised the stored values of the earlier seasons in the same way. This study uses the replay database.

\textbf{Spurious reporters.} Three callsign--grid pairs carried a distance that contradicted their own grid (Section~\ref{sec:qc}); this is the kind of entry that the cleaning step is meant to catch, and the live page now applies the same check.

\section{Conclusions}
\keyboxed{\begin{itemize}
\item The band is open (median at least 50 stations) from \busyFrom\ to \busyTo~UTC, with a median of about \busyMedAvg\ stations at midday; the long-distance windows lie at the edges of the day, not in its middle.
\item North America is reached most often in the morning (\peaknamerica\ UTC), Africa and Oceania in the afternoon and evening (\peakafrica\ UTC); the farthest reception was \recBest\km.
\item \pctDaysDXsix\,\% of days had a reporter beyond 6\,000\km\ with 200\,mW and a car-roof whip.
\item Within the day, the MUF matters more than anything else the logger records; K and A matter most for the longest paths; the daily solar flux has no detectable effect.
\item Midsummer is the weakest period (July 2024 and July 2026) and autumn the strongest (October 2025 and 2026), independently of the solar flux; for the same months, daytime reporting was higher in 2025 and 2026 than in 2024.
\end{itemize}}

\appendix
\section{Reproducing the analysis}
\label{app:repro}
The figures, tables and every number in the text are produced from the two SQLite databases by the script \texttt{doc/src/analyze.py} (Python~3 with NumPy and Matplotlib; the standard library for everything else). It writes \texttt{numbers.tex}, one \texttt{tab\_*.tex} file per table and the PDF figures to \texttt{doc/src/out}; the paper is then built with \texttt{pdflatex propagation\_study.tex} (twice). The databases themselves are not part of the repository.

\section*{Version history}
\begin{tabular}{@{}llp{9.5cm}@{}}\toprule Version & Date & Changes\\\midrule
1.0 & 6 November 2026 & First version: 2026 season only, 166 days.\\
1.1 & 6 November 2026 & Two earlier seasons (2024, 2025) rebuilt from \texttt{wspr.live}: \dateFrom\ to \dateTo, \nSnap\ snapshots; new section on the seasons; space-weather results explicitly limited to the live-logged days of 2026.\\
1.2 & 6 November 2026 & Measured the effect of late reports on the rebuilt seasons: live and rebuilt snapshots agree (ratio \cmpRatio).\\\bottomrule\end{tabular}

\begin{thebibliography}{99}\small
\bibitem{wsprnet} WSPRnet --- the Weak Signal Propagation Reporter Network. \url{https://wsprnet.org}
\bibitem{wsprlive} wspr.live --- ClickHouse interface to the WSPRnet spot database. \url{https://wspr.live}
\bibitem{wsjtx} J.~Taylor, K1JT, \emph{WSJT-X User Guide} (WSPR mode) and \emph{WSPR 2.0 User's Guide}. \url{https://wsjt.sourceforge.io}
\bibitem{zachtek} ZachTek, \emph{WSPR Desktop Transmitter} documentation. \url{https://www.zachtek.com/1012}
\bibitem{giro} Global Ionospheric Radio Observatory, Digital Ionogram Database (DIDBase). \url{https://giro.uml.edu/didbase/}
\bibitem{juliusruh} Leibniz Institute of Atmospheric Physics, Juliusruh Space Weather Monitor. \url{https://ionosonde.iap-kborn.de/actuellz.htm}
\bibitem{swpc} NOAA Space Weather Prediction Center, solar and geophysical indices. \url{https://www.swpc.noaa.gov}
\bibitem{hamqsl} N0NBH, solar-terrestrial data feed. \url{https://www.hamqsl.com/solar.html}
\end{thebibliography}
\end{document}
